% Generated draft: requires mathematical and source review.
\documentclass[11pt]{article}
\usepackage[margin=25mm]{geometry}
\usepackage{fontspec}
\setmainfont{Latin Modern Roman}
\usepackage{amsmath,amssymb,mathtools,unicode-math}
\setmathfont{Latin Modern Math}
\usepackage{enumitem,needspace,xurl,hyperref}
\setlength{\parindent}{0pt}
\setlength{\parskip}{5pt}
\setlength{\emergencystretch}{3em}
\newcommand{\N}{\mathbb{N}}
\newcommand{\Z}{\mathbb{Z}}
\newcommand{\Q}{\mathbb{Q}}
\newcommand{\R}{\mathbb{R}}
\newcommand{\C}{\mathbb{C}}
\newcommand{\F}{\mathbb{F}}
\newcommand{\A}{\mathbb{A}}
\newcommand{\PP}{\mathbb{P}}
\newcommand{\E}{\mathbb{E}}
\newcommand{\eps}{\varepsilon}
\newcommand{\dd}{\,\mathrm{d}}
\begin{document}
\section*{UnsolvedMath ID: 10600044}
\textbf{Source problem code:} catalogue entry
\subsection{Definitions and mathematical statement}
Let $\mathcal D_{\mathrm v}$ be virtual knot diagrams and let $\sim_{\mathrm v}$ be equivalence generated by the usual classical Reidemeister and virtual detour moves, with knot orientation conventions fixed consistently. An invariant $I:\mathcal D_{\mathrm v}\to S$ is complete if, for all $D,D'\in\mathcal D_{\mathrm v}$,
\[
 I(D)=I(D')\quad\Longleftrightarrow\quad D\sim_{\mathrm v}D'.
\]
Problem-list item 44 asks how to construct a complete invariant for virtual knots. The intended goal is a useful description of such an invariant and its target, beyond the tautological map sending each diagram to its own equivalence class. The excerpt does not require a particular algebraic target, algorithm, complexity bound, or effective equality test, so these must not be silently added as hypotheses.

The source motivates the request by inclusions of knot theories: it relates the embedding of classical knots into virtual knots to extending a quandle with peripheral structure, and gives the extension of Hurwitz's action for classical and virtual braids as an analogous example. These are supplied motivational assertions, not a definition of the requested virtual-knot invariant or a claim that the ordinary virtual knot group is complete.
\subsection{Short English statement}
Construct an explicitly described invariant that takes equal values exactly on equivalent virtual knots. The source does not specify the desired target or effectiveness requirements.
\subsection{Sources}
\begin{itemize}
\item Local registry: \url{data/unsolvedmath/problems.json}, permanent UnsolvedMath ID $10600044$, exact source code \texttt{catalogue entry}.
\item \emph{Unsolved Problems in Virtual Knot Theory and Combinatorial Knot Theory}, revised survey, problem-list item 44: \url{https://arxiv.org/abs/1409.2823}. Supplied access date: 2026-07-29; extraction: source-tex; compound item retained intact.
\item The registry names ``Quotient Quandles and the Fundamental Latin Alexander Quandle'' with incomplete identifier \texttt{arXiv:1404.xxxx} (2014). It also mentions incompleteness of virtual knot groups and low-crossing classifications, without precise supporting references.
\end{itemize}
\end{document}
